\documentclass[10pt]{beamer}
\usepackage[english]{babel}
\usepackage[utf8]{inputenc}
\usepackage[T1]{fontenc}
\usepackage{tikz}
\usepackage{changepage}
\usepackage{caption}
\usepackage{verbatim}
\usepackage{graphicx} % Required for inserting images
\usepackage{amsmath} % Mathematical symbols.
\usepackage{siunitx} % SI units.
\usepackage{hyperref} % Required for referring.
%\usepackage{minted}
\usepackage{listings}
\usepackage{xcolor}
\usepackage{tabularx}
\usepackage{booktabs}
\usepackage{array}
\usepackage{multirow}
%\usepackage{enumitem}
\usepackage[normalem]{ulem} % Per subratllar sense problemes de salts de línia
\usepackage{subcaption} % El paquet clau per a les subfigures
\usepackage{float}
\usepackage{gensymb}
\usetheme{metropolis}
% =========================
% PALETA DE COLORS PASTEL
% =========================
\definecolor{pastelSalmon}{HTML}{FFC09F} % #ffc09f
\definecolor{pastelYellow}{HTML}{FFEE93} % #ffee93
\definecolor{pastelRadioactiveYellow}{HTML}{FFDE20} % #ffde20
%\definecolor{pastelCream}{HTML}{FCF5C7}  % #fcf5c7
\definecolor{pastelCream}{HTML}{FFFBEA}  % #fffbea
\definecolor{pastelBlue}{HTML}{A0CED9}   % #a0ced9
\definecolor{pastelGreen}{HTML}{ADF7B6}  % #adf7b6
\definecolor{pastelPurple}{HTML}{D5AAFF} % #d5aaff
\definecolor{pastelPink}{HTML}{FFB5E8}   % #ffb5e8
\definecolor{pastelGray}{HTML}{CFCFC4}   % #cfcfc4
\definecolor{pastelOrange}{HTML}{FFAB73} % #ffab73
\definecolor{pastelMint}{HTML}{B9FBC0}   % #b9fbc0
\definecolor{pastelLavender}{HTML}{E2C6FF} % #e2c6ff
\definecolor{pastelCyan}{HTML}{9EF0F0}   % #9ef0f0
\definecolor{pastelDarkForest}{HTML}{009753} % #097534
\definecolor{pastelLightGray}{HTML}{E8E8E4} % #e8e8e4
\definecolor{pastelRed}{HTML}{d61010} % #d61010


% =========================
% CONFIGURACIÓ DE COLORS
% =========================
\setbeamercolor{background canvas}{bg=pastelCream}
\setbeamercolor{normal text}{fg=black,bg=pastelCream}
\setbeamercolor{structure}{fg=pastelRadioactiveYellow}
\setbeamercolor{frametitle}{bg=pastelRadioactiveYellow, fg=white}
\setbeamercolor{title}{fg=black,bg=pastelGreen}
\setbeamercolor{subtitle}{fg=black}
\setbeamercolor{progress bar in head/foot}{fg=pastelSalmon,bg=pastelCream}
\setbeamercolor{block title}{bg=pastelYellow, fg=black}
\setbeamercolor{block body}{bg=white, fg=black}

% =========================
% CONFIGURACIÓ DE FONTS
% =========================
\setbeamerfont{title}{size=\Large}
\setbeamerfont{subtitle}{size=\large}
\setbeamerfont{author}{size=\normalsize}
\setbeamerfont{institute}{size=\small}
\setbeamerfont{caption}{size=\footnotesize}

\setbeamersize{text margin left=1.5em, text margin right=1.5em}
% Cambiar "Cuadro" → "Tabla"
\addto\captionsspanish{%
  \renewcommand{\tablename}{Tabla}%
  \renewcommand{\figurename}{Figura}%
}
\makeatletter
\setlength{\metropolis@frametitle@padding}{0pt}
\setbeamertemplate{frametitle}{
  \nointerlineskip%
  \vskip\metropolis@frametitle@padding%
  \begin{beamercolorbox}[wd=\paperwidth,center]{frametitle}
    \vbox{}\vskip0ex%
    \strut\insertframetitle\strut\par
    \vskip0.75ex
  \end{beamercolorbox}
}
\makeatother


%=================================================================================================================================
%=================================================================================================================================


\title{TFM - ILC Vs. CILC $y$-maps and $\mu y$-maps}
%\subtitle{Presentación del treball}
\author{Joan Ribot Oliver}
\institute{Universidad de Cantabria}
\date{2026}


\begin{document}

\begin{frame}[plain]
    \titlepage
  % Logotip a baix a la dreta
  \begin{tikzpicture}[remember picture,overlay]
    \node[anchor=south east, xshift=-9.25cm, yshift=0.1cm] 
      at (current page.south east) {
\includegraphics[width=2.5cm]{/Users/joanribot/Desktop/Física/Màster UC - UIMP/2n semestre/TFM/Codes/Slides/logo_uc_800.jpg}};
  \end{tikzpicture}
\end{frame}


%=================================================================================================================================
% RESIDUAL COMPARISON: P-ILC VS P-CILC for both y and muy maps
%=================================================================================================================================

%------------------------------------------------------------------------------------------------------------------------------------
% mu = 2e-4, noise = 2 muK·arcmin, freqs = [30, 40, 90, 150, 220, 270] GHz
%------------------------------------------------------------------------------------------------------------------------------------

\begin{frame}[t]{$y$-maps Residual Comparison: P-ILC vs P-CILC}
  \centering
  4 layers: tSZ, $\mu$-distortion, CMB and noise ($0.2 \mu K \cdot$arcmin)\\
  $\mu = 2 \times 10^{-4}$, freqs = [30, 40, 90, 150, 220, 270] GHz
  \begin{figure}[H]
    \centering
    % Centratge absolut ignorant marges
    \begin{tikzpicture}[remember picture,overlay]
    \node[at=(current page.center), yshift=-0.4cm] {
      \includegraphics[width=\paperwidth]{/Users/joanribot/Desktop/Física/Màster UC - UIMP/2n semestre/TFM/Codes/Figures/ILC_VS_CILC_2e-04/y_maps/Residual_Comparison_P-CILC_VS_P-ILC.pdf}
    };
    \end{tikzpicture}
  \label{fig:y_maps_residual_comparison_2e-04}
  \end{figure}
\end{frame}

\begin{frame}[t]{$\mu y$-maps Residual Comparison: P-ILC vs P-CILC}
  \centering
  4 layers: tSZ, $\mu$-distortion, CMB and noise ($0.2 \mu K \cdot$arcmin)\\
  $\mu = 2 \times 10^{-4}$, freqs = [30, 40, 90, 150, 220, 270] GHz
  \begin{figure}[H]
    \centering
    % Centratge absolut ignorant marges
    \begin{tikzpicture}[remember picture,overlay]
    \node[at=(current page.center), yshift=-0.4cm] {
      \includegraphics[width=\paperwidth]{/Users/joanribot/Desktop/Física/Màster UC - UIMP/2n semestre/TFM/Codes/Figures/ILC_VS_CILC_2e-04/mu_y_maps/Residual_Comparison_P-CILC_VS_P-ILC.pdf}
    };
    \end{tikzpicture}
  \label{fig:mu_y_maps_residual_comparison_2e-04}
  \end{figure}
\end{frame}

%------------------------------------------------------------------------------------------------------------------------------------
% mu = 5e-4, noise = 2 muK·arcmin, freqs = [30, 40, 90, 150, 220, 270] GHz
%------------------------------------------------------------------------------------------------------------------------------------

\begin{frame}[t]{$y$-maps Residual Comparison: P-ILC vs P-CILC}
  \centering
  4 layers: tSZ, $\mu$-distortion, CMB and noise ($0.2 \mu K \cdot$arcmin)\\
  $\mu = 5 \times 10^{-4}$, freqs = [30, 40, 90, 150, 220, 270] GHz
  \begin{figure}[H]
    \centering
    % Centratge absolut ignorant marges
    \begin{tikzpicture}[remember picture,overlay]
    \node[at=(current page.center), yshift=-0.4cm] {
      \includegraphics[width=\paperwidth]{/Users/joanribot/Desktop/Física/Màster UC - UIMP/2n semestre/TFM/Codes/Figures/ILC_VS_CILC_5e-04/y_maps/Residual_Comparison_P-CILC_VS_P-ILC.pdf}
    };
    \end{tikzpicture}
  \label{fig:y_maps_residual_comparison_5e-04}
  \end{figure}
\end{frame}

\begin{frame}[t]{$\mu y$-maps Residual Comparison: P-ILC vs P-CILC}
  \centering
  4 layers: tSZ, $\mu$-distortion, CMB and noise ($0.2 \mu K \cdot$arcmin)\\
  $\mu = 5 \times 10^{-4}$, freqs = [30, 40, 90, 150, 220, 270] GHz
  \begin{figure}[H]
    \centering
    % Centratge absolut ignorant marges
    \begin{tikzpicture}[remember picture,overlay]
    \node[at=(current page.center), yshift=-0.4cm] {
      \includegraphics[width=\paperwidth]{/Users/joanribot/Desktop/Física/Màster UC - UIMP/2n semestre/TFM/Codes/Figures/ILC_VS_CILC_5e-04/mu_y_maps/Residual_Comparison_P-CILC_VS_P-ILC.pdf}
    };
    \end{tikzpicture}
  \label{fig:mu_y_maps_residual_comparison_5e-04}
  \end{figure}
\end{frame}

%------------------------------------------------------------------------------------------------------------------------------------
% mu = 2e-5, noise = 2 muK·arcmin, freqs = [30, 40, 90, 150, 220, 270] GHz
%------------------------------------------------------------------------------------------------------------------------------------

\begin{frame}[t]{$y$-maps Residual Comparison: P-ILC vs P-CILC}
  \centering
  4 layers: tSZ, $\mu$-distortion, CMB and noise ($0.2 \mu K \cdot$arcmin)\\
  $\mu = 2 \times 10^{-5}$, freqs = [30, 40, 90, 150, 220, 270] GHz
  \begin{figure}[H]
    \centering
    % Centratge absolut ignorant marges
    \begin{tikzpicture}[remember picture,overlay]
    \node[at=(current page.center), yshift=-0.4cm] {
      \includegraphics[width=\paperwidth]{/Users/joanribot/Desktop/Física/Màster UC - UIMP/2n semestre/TFM/Codes/Figures/ILC_VS_CILC_2e-05/y_maps/Residual_Comparison_P-CILC_VS_P-ILC.pdf}
    };
    \end{tikzpicture}
  \label{fig:y_maps_residual_comparison_2e-05}
  \end{figure}
\end{frame}

\begin{frame}[t]{$\mu y$-maps Residual Comparison: P-ILC vs P-CILC}
  \centering
  4 layers: tSZ, $\mu$-distortion, CMB and noise ($0.2 \mu K \cdot$arcmin)\\
  $\mu = 2 \times 10^{-5}$, freqs = [30, 40, 90, 150, 220, 270] GHz
  \begin{figure}[H]
    \centering
    % Centratge absolut ignorant marges
    \begin{tikzpicture}[remember picture,overlay]
    \node[at=(current page.center), yshift=-0.4cm] {
      \includegraphics[width=\paperwidth]{/Users/joanribot/Desktop/Física/Màster UC - UIMP/2n semestre/TFM/Codes/Figures/ILC_VS_CILC_2e-05/mu_y_maps/Residual_Comparison_P-CILC_VS_P-ILC.pdf}
    };
    \end{tikzpicture}
  \label{fig:mu_y_maps_residual_comparison_2e-05}
  \end{figure}
\end{frame}


%------------------------------------------------------------------------------------------------------------------------------------
% mu = 5e-5, noise = 2 muK·arcmin, freqs = [30, 40, 90, 150, 220, 270] GHz
%------------------------------------------------------------------------------------------------------------------------------------

\begin{frame}[t]{$y$-maps Residual Comparison: P-ILC vs P-CILC}
  \centering
  4 layers: tSZ, $\mu$-distortion, CMB and noise ($0.2 \mu K \cdot$arcmin)\\
  $\mu = 5 \times 10^{-5}$, freqs = [30, 40, 90, 150, 220, 270] GHz
  \begin{figure}[H]
    \centering
    % Centratge absolut ignorant marges
    \begin{tikzpicture}[remember picture,overlay]
    \node[at=(current page.center), yshift=-0.4cm] {
      \includegraphics[width=\paperwidth]{/Users/joanribot/Desktop/Física/Màster UC - UIMP/2n semestre/TFM/Codes/Figures/ILC_VS_CILC_5e-05/y_maps/Residual_Comparison_P-CILC_VS_P-ILC.pdf}
    };
    \end{tikzpicture}
  \label{fig:y_maps_residual_comparison_5e-05}
  \end{figure}
\end{frame}

\begin{frame}[t]{$\mu y$-maps Residual Comparison: P-ILC vs P-CILC}
  \centering
  4 layers: tSZ, $\mu$-distortion, CMB and noise ($0.2 \mu K \cdot$arcmin)\\
  $\mu = 5 \times 10^{-5}$, freqs = [30, 40, 90, 150, 220, 270] GHz
  \begin{figure}[H]
    \centering
    % Centratge absolut ignorant marges
    \begin{tikzpicture}[remember picture,overlay]
    \node[at=(current page.center), yshift=-0.4cm] {
      \includegraphics[width=\paperwidth]{/Users/joanribot/Desktop/Física/Màster UC - UIMP/2n semestre/TFM/Codes/Figures/ILC_VS_CILC_5e-05/mu_y_maps/Residual_Comparison_P-CILC_VS_P-ILC.pdf}
    };
    \end{tikzpicture}
  \label{fig:mu_y_maps_residual_comparison_5e-05}
  \end{figure}
\end{frame}



%=================================================================================================================================
% TT PLOTS: P-ILC vs P-CILC 
%=================================================================================================================================
\begin{comment}
%-------
% y maps
%-------
\begin{frame}[t]{TT plot for $y$ correlation: P-ILC Vs. P-CILC}
  \centering
  4 layers: tSZ, $\mu$-distortion, CMB and noise ($0.2 \mu K \cdot$arcmin)\\
  $\mu = 2 \times 10^{-4}$, freqs = [30, 40, 90, 150, 220, 270] GHz, $r$ = full map
  \begin{figure}[H]
    \centering
    \begin{tikzpicture}[remember picture,overlay]
    \node[at=(current page.center), yshift=-1.25cm] {
      \includegraphics[width=0.75\paperwidth]{/Users/joanribot/Desktop/Física/Màster UC - UIMP/2n semestre/TFM/Codes/Figures/ILC_VS_CILC_2e-04/TT_plots/TT_Plot_y_ILC_VS_CILC.pdf}
    };
    \end{tikzpicture}
  \label{fig:PCILC_zoomed_y_maps_0.2muK_2e-04}
  \end{figure}
\end{frame}

\begin{frame}[t]{TT plot for $y$ correlation: P-ILC Vs. P-CILC}
  \centering
  4 layers: tSZ, $\mu$-distortion, CMB and noise ($0.2 \mu K \cdot$arcmin)\\
  $\mu = 2 \times 10^{-4}$, freqs = [30, 40, 90, 150, 220, 270] GHz, $r$ = 10 pixels
  \begin{figure}[H]
    \centering
    \begin{tikzpicture}[remember picture,overlay]
    \node[at=(current page.center), yshift=-1.25cm] {
      \includegraphics[width=0.75\paperwidth]{/Users/joanribot/Desktop/Física/Màster UC - UIMP/2n semestre/TFM/Codes/Figures/ILC_VS_CILC_2e-04/TT_plots/TT_Plot_y_ILC_VS_CILC_10p_radius.pdf}
    };
    \end{tikzpicture}
  \label{fig:PCILC_zoomed_y_maps_0.2muK_10p}
  \end{figure}
\end{frame}

%----------
% mu y maps
%----------
\begin{frame}[t]{TT plot for $\mu y$ correlation: P-ILC Vs. P-CILC}
  \centering
  4 layers: tSZ, $\mu$-distortion, CMB and noise ($0.2 \mu K \cdot$arcmin)\\
  $\mu = 2 \times 10^{-4}$, freqs = [30, 40, 90, 150, 220, 270] GHz, $r$ = full map
  \begin{figure}[H]
    \centering
    \begin{tikzpicture}[remember picture,overlay]
    \node[at=(current page.center), yshift=-1.25cm] {
      \includegraphics[width=0.75\paperwidth]{/Users/joanribot/Desktop/Física/Màster UC - UIMP/2n semestre/TFM/Codes/Figures/ILC_VS_CILC_2e-04/TT_plots/TT_Plot_muy_ILC_VS_CILC.pdf}
    };
    \end{tikzpicture}
  \label{fig:PCILC_zoomed_muy_maps_0.2muK_2e-04}
  \end{figure}
\end{frame}

\begin{frame}[t]{TT plot for $\mu y$ correlation: P-ILC Vs. P-CILC}
  \centering
  4 layers: tSZ, $\mu$-distortion, CMB and noise ($0.2 \mu K \cdot$arcmin)\\
  $\mu = 2 \times 10^{-4}$, freqs = [30, 40, 90, 150, 220, 270] GHz, $r$ = 10 pixels
  \begin{figure}[H]
    \centering
    \begin{tikzpicture}[remember picture,overlay]
    \node[at=(current page.center), yshift=-1.25cm] {
      \includegraphics[width=0.75\paperwidth]{/Users/joanribot/Desktop/Física/Màster UC - UIMP/2n semestre/TFM/Codes/Figures/ILC_VS_CILC_2e-04/TT_plots/TT_Plot_muy_ILC_VS_CILC_10p_radius.pdf}
    };
    \end{tikzpicture}
  \label{fig:PCILC_zoomed_muy_maps_0.2muK_10p}
  \end{figure}
\end{frame}
\end{comment}



\begin{comment}


%\begin{frame}[plain]
%  % Fons d’imatge
%  \begin{tikzpicture}[remember picture,overlay]
%    \node[at=(current page.center)] {
%      \includegraphics[width=\paperwidth,height=\paperheight]{UGC_2885.tif.jpg}
%    };
%  \end{tikzpicture}
%  \begin{tikzpicture}[remember picture,overlay]
%    \node[at=(current page.center), yshift=3cm] {
%      \Huge\color{white} Tanda de preguntas
%    };
%  \end{tikzpicture}
%\end{frame}


%\begin{frame}{}
%    \begin{tikzpicture}[remember picture,overlay]
%    \node[at=(current page.center)] {
%      \includegraphics[width=\paperwidth,height=\paperheight]{Super_tunning_fork.jpg}
%    };
%  \end{tikzpicture}
%\end{frame}


%\begin{frame}[plain]
%  \begin{tikzpicture}[remember picture,overlay]
%    \node[at=(current page.center)] {
%      \includegraphics[width=\paperwidth,height=\paperheight]{UGC_2885_VR.png}
%    };
%  \end{tikzpicture}
%\end{frame}

%\begin{frame}[plain] % "plain" opcional per eliminar marges del tema
%  \begin{tikzpicture}[remember picture,overlay]
%    % (x,y) es una posición relativa a toda la diapositiva
%    \node[anchor=north west, xshift=1cm, yshift=-1.5cm] 
%      at (current page.north west) {
%        \parbox{0.9\textwidth}{
%          \textbf{Resultados:} \\
%          En esta diapositiva colocamos el texto justo donde queremos.
%          \begin{itemize}
%            \item Corresponde a una posición específica.
%            \item No se moverá aunque cambies el resto del contenido.
%          \end{itemize}
%        }
%      };
%  \end{tikzpicture}
%\end{frame}
\end{comment}

\end{document}
